\section{Bharadiya, Thomas e Ahmed (2023) -- Rise of Artificial Intelligence in Business and Industry}

\subsection*{Resumo}

O artigo é uma revisão, classificada pela revista como \emph{Review Article} (p.~85). Segundo os autores, revisões anteriores sobre IA em negócios se concentram em aplicações específicas, como \emph{business intelligence} e sistemas de apoio à decisão, ou em domínios específicos (p.~86). O objetivo declarado é oferecer uma visão geral da pesquisa sobre IA em negócios, analisando sua história e o estado da arte e projetando tendências futuras (p.~86--87).

O texto começa definindo inteligência artificial, \emph{machine learning} e \emph{deep learning} (p.~87). Segundo os autores, a IA complementa, e não substitui, a inteligência humana e ajuda a prever os resultados de cada ação e a simplificar a tomada de decisão (p.~88).

Em seguida, os autores descrevem aplicações da IA em negócios, como cibersegurança, gestão de clientes, segmentação de público, detecção de fraude e cadeia de suprimentos (p.~88--90), e em marketing, com base na classificação de Davenport e Ronanki (p.~90). A maior parte do texto percorre aplicações por setor, entre eles saúde, varejo, alimentos, finanças, transporte e energia (p.~92--97). No setor financeiro, são citados a análise automatizada de pedidos de crédito, os \emph{robo-advisors}, os \emph{chatbots} em seguros e a detecção de fraude(p.~94).

A conclusão afirma que a IA está promovendo uma revolução nos negócios, na economia e na sociedade (p.~97).

\subsection*{Natureza do texto}

Revisão narrativa, sem descrição de como os estudos foram selecionados. Na leitura, algumas citações numeradas não correspondiam à lista de referências (p.~86 e~98). Serve como panorama do uso de IA em negócios; afirmações específicas devem ser verificadas nas fontes originais.

\subsection*{O que o texto trouxe de novo}

O principal ponto que o texto trouxe foi que o valor da IA para uma empresa não está só no modelo ou no algoritmo, mas em transformar dados em informação ou ação que melhora um processo, apoia uma decisão ou reduz algum problema da organização. Isso aparece quando os autores afirmam que a IA ajuda a prever os resultados de cada ação e a simplificar a tomada de decisão (p.~88).

Também ficou claro que a mesma base tecnológica pode assumir funções diferentes conforme o contexto. Ao longo do texto, \emph{machine learning} aparece aplicado à detecção de fraude, à recomendação de produtos, à manutenção preditiva e à análise de clientes (p.~87--97). Isso ajuda a distinguir duas coisas: desenvolver um modelo de IA e usá-lo para gerar valor em uma aplicação real.

\subsection*{O que extrair para um trabalho autoral}

O texto serve como contexto para a introdução do projeto: mostra que a IA vem sendo aplicada em vários setores, inclusive no financeiro. Não deve ser usado como base para afirmações específicas.

A partir dele, uma estrutura possível para o projeto é: dados $\rightarrow$ modelo $\rightarrow$ informação $\rightarrow$ decisão $\rightarrow$ resultado. Essa estrutura é uma interpretação minha, apoiada na ideia dos autores de que a IA ajuda a prever resultados e a simplificar decisões (p.~88).

Há também uma lacuna: as aplicações financeiras citadas (crédito, \emph{robo-advisors}, seguros e fraude) não incluem decisões de investimento com avaliação de resultado (p.~94). É nesse espaço que o projeto se situa.

\subsection*{Autores para catalogar}

Citados pelos autores, ainda não lidos (status C). Como a numeração das citações não é confiável, a busca deve ser feita por nome e ano.

\begin{itemize}
  \item Duan, Edwards e Dwivedi (2019) -- IA e
        decisão; relevância alta.
  \item Jarrahi (2018) -- humano e IA na decisão
        organizacional; relevância alta.
  \item Moro, Cortez e Rita (2015) -- \emph{business
        intelligence} em bancos; relevância média.
  \item Riikkinen et~al.~(2018) -- IA em seguros;
        relevância média.
\end{itemize}

Os demais autores citados já constam no Catálogo, com relevância baixa para o problema do projeto.